\chapter*{Заключение}
\addcontentsline{toc}{chapter}{Заключение}
\label{ch:conclusion}

В ходе расчётно-графической работы спроектирован и реализован тестовый
веб-видеохостинг на современном стеке веб-технологий.

Выполнены все поставленные задачи:
\begin{itemize}
  \item спроектирован в Figma макет интерфейса, его дизайн-система и
        состав экранов перенесены в код;
  \item разработана серверная часть на Django и Django REST Framework:
        модель данных, авторизация по паролю и гостевые аккаунты с
        выдачей JWT-токенов, загрузка и потоковая отдача видео с
        поддержкой Range;
  \item разработано клиентское приложение на React с маршрутизацией,
        контекстом авторизации и компонентами дизайн-системы;
  \item серверная и клиентская части покрыты автоматическими тестами;
  \item платформа развёрнута на собственном сервере за обратным прокси
        Caddy с доступом по HTTPS, документация API оформлена в Git.
\end{itemize}

В процессе работы получены практические навыки построения
клиент-серверного веб-приложения: проектирование REST-API, реализация
авторизации на JWT-токенах, разработка одностраничного приложения на
React и контейнерное развёртывание.

Кодовая база организована по принципам модульности, DRY и единственной
ответственности, снабжена тестами и документацией, поэтому пригодна для
дальнейшего расширения. Возможные направления развития: совместный
просмотр видео в общих комнатах, поиск и рекомендации, транскодирование
загружаемых видео в несколько качеств с адаптивным стримингом.

\endinput
